% =====================================================================
% Figure contract
% 核心结论：§5 数据体系沿监督粒度由粗到细递升，同时来源轴与各层正交。
% 图原型：中等档 / 定制横向阶梯（同批图配色与 preamble 风格）。
% 证据层级：hero = 五级监督粒度阶梯；支撑 = 右侧双分支与底部来源轴。
% 能不能更少？：不能；五级阶梯、两张支卡、三类来源与两条小注均为指定内容。
%
% Step ① 设计文档（Form A / 中等复杂度）
% 1 领域：Computer-Use Agents 中文综述；技术术语保留英文。
% 2 格式：TikZ standalone；XeLaTeX + fontspec + xeCJK；全图无旋转中文。
% 3 画布：W_print = 22 cm；目标最小字号 7.5 pt；W_ink <= 22 cm。
% 4 布局：五张主卡左→右逐级微升；第 5 卡向右分出上下两张支卡。
% 5 连线：主阶梯为短实线箭头；支卡为虚线箭头；来源轴用 3 条浅色虚线 leader。
% 6 空间：支卡占右侧独立列；来源轴横贯底部；两条灰注分居轴下左右。
% 7 强调：五卡同一蓝灰色系逐级加深；支卡与来源标签降低视觉重量。
% 8 密度：中等；无数值，不嵌入无依据的图表或公式。
% 9 ASCII 草图：
%       [GUI] ↗ [Grounding] ↗ [Atomic] ↗ [Trajectory] ↗ [Traces] - - > [失败·偏好·安全]
%                                                               \ - - > [个性化]
%       ··················· 来源轴：Human / Synthetic / Agent-generated ···················
%       [左下小注]                                                     [右下小注]
% 10 可视化决策：输入为分类关系而非数值，故只用结构、色阶与正交轴表达。
% Pre-flight：positioning 相对布局；连线均接 node anchors；直线无 rounded corners；
%              leader 只落在卡片下边界，不穿框；文字用固定 text width + 手动换行。
% =====================================================================

\documentclass[border=8pt]{standalone}
\usepackage{fontspec}
\usepackage{xeCJK}
\usepackage{xcolor}
\usepackage{tikz}
\setCJKmainfont{Heiti SC}
\setCJKsansfont{Heiti SC}
\usetikzlibrary{arrows.meta,bending,positioning,calc}

% Low-saturation palette aligned with cua-survey-fig1-overview.tex.
\definecolor{ink}{HTML}{27364A}
\definecolor{muted}{HTML}{748092}
\definecolor{mainLine}{HTML}{4F667C}
\definecolor{stepLineOne}{HTML}{8DA0AD}
\definecolor{stepLineTwo}{HTML}{8095A4}
\definecolor{stepLineThree}{HTML}{738998}
\definecolor{stepLineFour}{HTML}{667D8D}
\definecolor{stepLineFive}{HTML}{566F81}
\definecolor{stepFillOne}{HTML}{EFF4F6}
\definecolor{stepFillTwo}{HTML}{E8F0F3}
\definecolor{stepFillThree}{HTML}{E0EAEE}
\definecolor{stepFillFour}{HTML}{D7E4E9}
\definecolor{stepFillFive}{HTML}{CDDEE5}
\definecolor{purpleLine}{HTML}{817594}
\definecolor{purpleFill}{HTML}{EEEAF2}
\definecolor{tealLine}{HTML}{648D8A}
\definecolor{tealFill}{HTML}{E5F0EE}
\definecolor{roseLine}{HTML}{9A777B}
\definecolor{roseFill}{HTML}{F3E9EA}
\definecolor{axisFill}{HTML}{F1F4F6}
\definecolor{axisLine}{HTML}{A7B1BB}

\tikzset{
  every node/.style={
    font=\sffamily\fontsize{8.6}{10.4}\selectfont,
    text=ink
  },
  stepcard/.style={
    rectangle,
    rounded corners=4pt,
    line width=0.8pt,
    minimum width=2.94cm,
    minimum height=1.34cm,
    text width=2.52cm,
    align=center,
    inner xsep=6pt,
    inner ysep=5pt,
    font=\sffamily\fontsize{8.7}{10.5}\selectfont\bfseries
  },
  branchcard/.style={
    rectangle,
    rounded corners=4pt,
    line width=0.75pt,
    minimum width=3.72cm,
    minimum height=1.12cm,
    text width=3.30cm,
    align=center,
    inner xsep=6pt,
    inner ysep=4.5pt,
    font=\sffamily\fontsize{8.4}{10.2}\selectfont\bfseries
  },
  sourcebar/.style={
    rectangle,
    rounded corners=4pt,
    draw=axisLine,
    fill=axisFill,
    line width=0.6pt,
    minimum width=20.30cm,
    minimum height=1.06cm,
    inner sep=0pt
  },
  sourcetag/.style={
    rectangle,
    rounded corners=3pt,
    line width=0.55pt,
    minimum height=0.55cm,
    align=center,
    inner xsep=7pt,
    inner ysep=3pt,
    font=\sffamily\fontsize{8.0}{9.6}\selectfont\bfseries
  },
  title/.style={
    font=\sffamily\fontsize{13.2}{15.8}\selectfont\bfseries,
    text=mainLine
  },
  note/.style={
    font=\sffamily\fontsize{7.6}{9.3}\selectfont,
    text=muted,
    inner sep=0pt
  },
  arrow short/.style={
    -{Stealth[length=3pt 1.0,width'=0pt 0.5,sep=0pt 0.3]},
    line width=0.8pt,
    shorten >=1pt,
    shorten <=0pt,
    color=mainLine
  },
  branch arrow/.style={
    -{Stealth[length=3.6pt 1.0,width'=0pt 0.55,sep=0pt 0.3]},
    line width=0.75pt,
    densely dashed,
    shorten >=1.5pt,
    shorten <=0.5pt,
    color=purpleLine
  },
  source leader/.style={
    densely dashed,
    line width=0.45pt,
    color=axisLine!78,
    shorten >=1.5pt,
    shorten <=0pt
  }
}

\begin{document}
\sffamily
\begin{tikzpicture}

% Five-step supervision-granularity staircase, positioned by construction.
\node[stepcard,fill=stepFillOne,draw=stepLineOne] (s1) at (0,0)
  {GUI 理解 §5.1};
\node[stepcard,fill=stepFillTwo,draw=stepLineTwo,
      right=0.32cm of s1,yshift=0.34cm] (s2)
  {Grounding §5.2};
\node[stepcard,fill=stepFillThree,draw=stepLineThree,
      right=0.32cm of s2,yshift=0.34cm] (s3)
  {Atomic action\\§5.3};
\node[stepcard,fill=stepFillFour,draw=stepLineFour,
      right=0.32cm of s3,yshift=0.34cm] (s4)
  {Trajectory §5.4};
\node[stepcard,fill=stepFillFive,draw=stepLineFive,
      right=0.32cm of s4,yshift=0.34cm] (s5)
  {规划推理 traces\\§5.5};

% Short, straight connectors: no rounded corners on single-segment paths.
\draw[arrow short] (s1.east) -- (s2.west);
\draw[arrow short] (s2.east) -- (s3.west);
\draw[arrow short] (s3.east) -- (s4.west);
\draw[arrow short] (s4.east) -- (s5.west);

% Two quiet branches from planning/reasoning traces.
\node[branchcard,fill=roseFill,draw=roseLine,
      right=0.55cm of s5.east,anchor=west,yshift=0.88cm] (b1)
  {失败·偏好·安全\\§5.6–5.7};
\node[branchcard,fill=tealFill,draw=tealLine,
      right=0.55cm of s5.east,anchor=west,yshift=-0.70cm] (b2)
  {个性化 §5.8};
\draw[branch arrow] (s5.north east) -- (b1.west);
\draw[branch arrow] (s5.south east) -- (b2.west);

% Figure title, centered over the full staircase-plus-branch span.
\coordinate (horizontalCenter) at ($(s1.west)!0.5!(b1.east)$);
\coordinate (titleHeight) at ([yshift=0.56cm]b1.north);
\node[title,anchor=south] at (horizontalCenter |- titleHeight)
  {数据分层：监督粒度阶梯 × 来源轴（§5）};

% Orthogonal source axis across the bottom.
\node[sourcebar,anchor=north west] (source)
  at ([yshift=-1.15cm]s1.south west) {};

% Three sparse leaders terminate at card borders and never enter a frame.
\draw[source leader] (source.north -| s1.south) -- (s1.south);
\draw[source leader] (source.north -| s3.south) -- (s3.south);
\draw[source leader] (source.north -| s5.south) -- (s5.south);

\node[anchor=west,font=\sffamily\fontsize{8.7}{10.4}\selectfont\bfseries,
      text=mainLine] (sourceLabel) at ([xshift=0.36cm]source.west)
  {来源轴 §5.9};
\node[sourcetag,fill=stepFillTwo,draw=stepLineTwo,
      right=0.74cm of sourceLabel] (human) {Human};
\node[sourcetag,fill=purpleFill,draw=purpleLine,
      right=0.54cm of human] (synthetic) {Synthetic};
\node[sourcetag,fill=tealFill,draw=tealLine,
      right=0.54cm of synthetic] (agentgen) {Agent-generated};

% Gray footnotes: processing/quality at lower left, transfer gap at lower right.
\node[note,anchor=north west,text width=6.8cm,align=left]
  at ([yshift=-0.34cm]source.south west)
  {处理与混配 §5.10；质量·隐私·污染 §5.11};
\node[note,anchor=north east,text width=11.8cm,align=right]
  at ([yshift=-0.34cm]source.south east)
  {理解层高分不自动传导为端到端能力——理解与行动之间存在断层（§5.1）};

\end{tikzpicture}
\end{document}
